\scanpage{200}
\mspage{199}

ALGÈBRES de LIE SEMI-SIMPLES (5) \hfill Tribu n° \unsure{57|59}

\uline{n° actuellement 4} -- Dire ``produit fini'' dans la déf. 3. Dire que $\operatorname{Autel}(g)$ est invariant dans $\operatorname{Aut}(g)$. Au milieu de la p. 12 raisonner sur $\underline{g}/\underline{h}$, (au lieu de Fitting).

Ajouter `` $(t\in k)$ '' à la l.12 p. 13 ; dire ``une seule valeur'' à la l.14. Il s'agit d'une inclusion à la l.18. Dire `` image par $F$ '' à la l. 19.

Dans l'énoncé du th. 2 dire ``le groupe $\operatorname{Autel}(G)$'' opère transitivement'' sur les Cartans. Dire ``la conclusion de (iii)'' dans la remarque p. 14.

Ajouter le cor. suivant : soient $\underline{h}\subset\underline{g}'\subset\underline{g}$ ; si $\underline{g}'$ a même rang que $\underline{g}$ et contient un élément régulier de $\underline{g}$, alors $\underline{h}$ est une Cartan de $\underline{g}$ si et seult. si c'est une Cartan de $\underline{g}'$, (utiliser la prop. ajoutée au n° 3).

\uline{n° 6} -- Dire partout ``commutative'', (au lieu de ``abélienne''). Ajouter ``car nilpotente'' à la l.9, p. 16. Utiliser la prop. ajoutée au n° 3 dans la prop. 14.

\uline{n° 7} -- Style du th. 4 : ``le groupe des... opère transitivement...''. Idem dans le cor. p. 20. Dire ``composante neutre'' à la p. 21, l.3. Ajouter que, si $\underline{g}$ est une alg. Lie sur $K$, et si $k$ est un sous-corps de $K$, les Cartans de $\underline{g}$ considérée comme algèbre sur $k$ sont les mêmes que celles de $\underline{g}$ considérée comme algèbre sur $K$.

§2 -- \uuline{Les représentations de} $\mathrm{sl}_2$.

Eviter les points, (p. 24). Dire : ``sous-algèbre diagonale''.

\uline{n° 1} -- Dans la remarque, (p. 25), dire ``(car $X$ est un crochet)''. Généraliser la déf. des poids à tous les vecteurs propres de $H$. Un exer. montrera les séries génératrices correspondant à (2). A la p. 25, l.2 bas, dire `` d'où, en divisant par $n!$ ''.
